\documentclass[11pt]{article}

\usepackage[a4paper,margin=1in]{geometry}
\usepackage[T1]{fontenc}
\usepackage[utf8]{inputenc}
\usepackage{lmodern}
\usepackage{microtype}
\usepackage{amsmath,amssymb,amsfonts}
\usepackage{booktabs}
\usepackage{enumitem}
\usepackage{xcolor}
\usepackage{hyperref}

\emergencystretch=3em

\hypersetup{
  colorlinks=true,
  linkcolor=blue!55!black,
  citecolor=blue!55!black,
  urlcolor=blue!55!black,
  pdftitle={LMP Layer Documentation},
  pdfauthor={LMP project}
}

\newcommand{\LMP}{\texttt{LMP}}
\newcommand{\code}[1]{\texttt{\detokenize{#1}}}
\newcommand{\field}{\psi}
\newcommand{\site}{n}
\newcommand{\muindex}{\mu}
\newcommand{\U}{U}
\newcommand{\gauge}{g}
\newcommand{\C}{\mathbb{C}}
\newcommand{\Tr}{\operatorname{Tr}}
\newcommand{\diag}{\operatorname{diag}}
\newcommand{\norm}[1]{\left\lVert #1 \right\rVert}
\newcommand{\inner}[2]{\left\langle #1, #2 \right\rangle}
\newcommand{\pathop}{\mathcal{T}}
\newcommand{\paths}{\mathcal{P}}
\newcommand{\spin}{\mathcal{S}}
\newcommand{\dirac}{D}
\newcommand{\dd}{\mathrm{d}}

\title{\LMP{} Layer Documentation\\
\large Gauge-Covariant Machine-Learning Layers for Lattice QCD}
\author{LMP project}
\date{\today}

\begin{document}
\maketitle

\begin{abstract}
\LMP{} is a Julia package for building machine-learning models directly on
lattice QCD (LQCD) fermion fields. The package is organized around a simple
principle: learned operations should be compatible with local gauge covariance,
rather than treating lattice fields as ordinary dense arrays. Its main layers
therefore combine gauge-link parallel transport, spinor-space trainable
matrices, attention-style contractions, and Rayleigh--Ritz objectives for
low-mode learning. This document presents the scientific motivation, the
mathematical forward maps, and the implementation locations of the principal
layers. It is intended as a compact technical overview for collaborators who
know LQCD and want to understand what is being learned by the current \LMP{}
architecture.
\end{abstract}

\section*{Executive summary}

\LMP{} can be viewed as a physics-informed neural layer library for Wilson
fermions. Its design differs from a standard neural-network package in three
ways:
\begin{enumerate}[leftmargin=*]
  \item \textbf{Gauge covariance is built into the feature construction.}
        Neighboring fermion fields are not mixed directly; they are first
        parallel-transported with the gauge links to a common site.
  \item \textbf{Trainable parameters act primarily in spinor and channel
        space.} This keeps the color gauge transformation law intact while
        allowing expressive learned maps.
  \item \textbf{Several layers optimize spectral information.} Rayleigh
        quotients, Ritz values, orthogonalization, and blockwise amplitudes
        target low-mode subspaces relevant to deflation and multigrid-style
        solvers.
\end{enumerate}

The result is a collection of layers that are interpretable in LQCD language:
parallel-transport convolutions, gauge-covariant attention mechanisms, and
variational objectives for Dirac-operator spectra.

\tableofcontents
\newpage

\section{Purpose and high-level architecture}

\LMP{} treats a neural-network layer as an operator acting on Wilson fermion
fields, vectors of fermion fields, or scalar objective values built from such
fields. The central design idea is to replace ordinary image-like convolution
with \emph{parallel-transport convolution}: nearby lattice information is first
transported along gauge links to a common site, and only then mixed by trainable
spin-space matrices. This makes the construction naturally compatible with the
local gauge symmetry of LQCD.

This document is written from the viewpoint of a researcher inspecting the
architecture, not from the viewpoint of an API reference. Each layer is
therefore described by:
\begin{itemize}[leftmargin=*]
  \item the physical object it manipulates;
  \item the mathematical forward map implemented by the layer;
  \item the reason this construction is natural for LQCD;
  \item the source files where the implementation can be checked.
\end{itemize}

\subsection{Core design pattern}

Most field-valued \LMP{} layers follow the same structural pattern:
\[
  \boxed{
  \text{gauge transport}
  \quad\longrightarrow\quad
  \text{spin/channel mixing}
  \quad\longrightarrow\quad
  \text{optional spectral objective}
  }.
\]
The first step ensures that all features transform in the same local color
frame. The second step supplies learnable expressivity without violating that
transformation law. The third step, when present, connects the learned
representation to Dirac-operator low modes.

\subsection{Current implementation map}

The source layout relevant for layer documentation is:
\begin{itemize}[leftmargin=*]
  \item \code{src/LMP.jl}: top-level module, includes submodules, and exports
        public constructors.
  \item \code{src/Layer_types.jl}: layer structs and activation-function types.
  \item \code{src/Layer_sequence.jl}: forward maps, normalization, objective
        layers, sequence execution, CPU/GPU conversion helpers, and many
        gradient/update definitions.
  \item \code{src/elementry_base.jl}: fermion-field algebra, spin matrices,
        path data structures, hopping/parallel transport, inner products, block
        utilities, and clover-field utilities.
  \item \code{src/Gauge.jl}: gauge-transformation helpers.
  \item \code{src/gpu.jl}: GPU versions of spin matrices, packing routines,
        batched kernels, and GPU helper functions.
\end{itemize}

\subsection{What to focus on when presenting \LMP{}}

For a short introduction to collaborators, the most important points are:
\begin{itemize}[leftmargin=*]
  \item \code{Hopping_operator} is the geometric backbone: it constructs
        gauge-covariant path features.
  \item \code{PTC_Layer}, \code{CPTC_Layer}, and \code{PTLayer} are the main
        convolution-like layers.
  \item The attention layers replace scalar image attention with contractions
        of covariantly transported spinor fields.
  \item \code{RayleighQuotient_Layer} and \code{Ritz_Layer} turn the model into
        a variational tool for low-mode subspace learning.
  \item \code{GramSchmidt_Layer}, \code{Orthogonal_Layer}, and
        \code{BlockScalar_Layer} provide numerical structure needed by
        multi-vector and domain-decomposed representations.
\end{itemize}

\section{Notation and LQCD background}

\subsection{Lattice fields}

Let the four-dimensional lattice be
\[
  \Lambda =
  \{(n_1,n_2,n_3,n_4) \mid 1 \le n_\mu \le N_\mu,\ \mu=1,\dots,4\}.
\]
The gauge links are matrices
\[
  \U_\mu(n) \in SU(3),
\]
living on links from site $n$ to site $n+\hat\mu$. A Wilson fermion field is
represented as
\[
  \field(n) \in \C^{N_c} \otimes \C^{N_s},
  \qquad N_c=3,\quad N_s=4,
\]
so each site carries color and spin components.

In the implementation, Wilson fermion fields are represented by
\code{WF3x1CPU} or \code{WF3x1GPU}, type aliases for Wilson-fermion objects
from \code{LatticeDiracOperators}. Trainable spin matrices are represented by
\code{Spin_matrix} on CPU and \code{Spin_matrix_GPU} on GPU.

\begin{table}[h]
\centering
\small
\begin{tabular}{@{}ll@{}}
\toprule
Symbol & Meaning \\
\midrule
$n\in\Lambda$ & Lattice site \\
$\mu=1,\dots,4$ & Lattice direction \\
$\U_\mu(n)$ & Gauge link from $n$ to $n+\hat\mu$ \\
$\field(n)$ & Wilson fermion field with color and spin indices \\
$W$ & Trainable $4\times4$ spin matrix \\
$\pathop_p[\U]$ & Gauge-link transport along path $p$ \\
$\inner{x}{y}$ & Lattice color-spin inner product \\
$\dirac$ & Dirac-type operator used in spectral objectives \\
\bottomrule
\end{tabular}
\caption{Notation used throughout the document.}
\end{table}

\subsection{Gauge transformations}

A local gauge transformation is a site-wise $SU(3)$ matrix $\gauge(n)$. The
fermion and gauge fields transform as
\begin{align}
  \field(n) &\mapsto \field^g(n) = \gauge(n)\field(n),\\
  \U_\mu(n) &\mapsto \U_\mu^g(n)
    = \gauge(n)\U_\mu(n)\gauge^\dagger(n+\hat\mu).
\end{align}
This exact link rule is implemented in \code{gUgshift!} in \code{src/Gauge.jl}.
The function multiplies $\gauge(n)\U_\mu(n)$ and then right-multiplies by the
shifted gauge transformation $\gauge^\dagger(n+\hat\mu)$.

\subsection{Fermion inner product and norm}

The package uses the lattice field inner product implemented through
\code{dot(x,f)}:
\[
  \inner{x}{y}
  =
  \sum_{n\in\Lambda}
  \sum_{c=1}^{3}
  \sum_{\alpha=1}^{4}
  x_{c,\alpha}(n)^{*} y_{c,\alpha}(n).
\]
The squared norm is
\[
  \norm{x}^{2} = \inner{x}{x}.
\]
These appear as \code{inner_product(x,f)} and \code{norm2(x)} in
\code{src/elementry_base.jl}.

\section{Gauge-covariant paths and parallel transport}

\subsection{Path representation}

\LMP{} uses integer directions to represent lattice paths. A direction
$+\mu$ or $-\mu$ means a one-link displacement in the positive or negative
lattice direction. The special direction $0$ represents the identity path.
The relevant objects are:
\begin{itemize}[leftmargin=*]
  \item \code{Path}: one ordered list of steps.
  \item \code{Paths}: a collection of paths.
  \item \code{Onestep()}: returns the identity path and all positive/negative
        one-hop directions.
  \item \code{reverse_paths}: reverses a path and flips all signs, giving the
        adjoint transport path.
\end{itemize}

\subsection{Parallel transport operator}

For a path $p=(s_1,\dots,s_L)$ with $s_i\in\{0,\pm1,\dots,\pm4\}$, define the
path-transported field
\[
  \pathop_p[\U]\field(n)
  =
  \left(\prod_{\ell\in p}^{\mathrm{path}}
  \U_{\ell}^{(\pm)}(\cdot)\right)
  \field(n+\Delta p),
\]
where $\Delta p$ is the endpoint displacement and the product uses either a
link or its Hermitian adjoint depending on path orientation. In the code this is
implemented by \code{Hopping_operator}.

The CPU implementation in \code{src/elementry_base.jl} follows the rule:
\begin{align}
  s=+\mu:\quad
  \field'(n)
  &=
  \U_\mu^\dagger(n-\hat\mu)\,
  \field(n-\hat\mu),
  \\
  s=-\mu:\quad
  \field'(n)
  &=
  \U_\mu(n)\,
  \field(n+\hat\mu).
\end{align}
Boundary conditions are applied during the shift. The important property is
gauge covariance:
\[
  \pathop_p[\U^g]\field^g(n)
  =
  \gauge(n)\,\pathop_p[\U]\field(n).
\]
Therefore, all path-transported features in a layer transform at the same site
as the original field, making local spin-space mixing gauge compatible.

\section{Spin matrices and Clifford basis}

\subsection{Spin-space trainable weights}

\code{Spin_matrix} wraps a $4\times4$ complex matrix $W\in\C^{4\times4}$. It
acts on the spin index of a Wilson fermion field at each color and lattice site:
\[
  (W\field)_{c,\alpha}(n)
  =
  \sum_{\beta=1}^{4} W_{\alpha\beta}\field_{c,\beta}(n).
\]
This action does not mix color and therefore commutes with local color-gauge
transformations:
\[
  W(\gauge\field)=\gauge(W\field).
\]
That commutation is the reason spin matrices can be learned without breaking
gauge covariance.

\subsection{Clifford-transformed weights}

\code{CPTC_Layer} and several attention layers use \code{Clifford(W)}. The
source defines a set of gamma/Clifford matrices and maps raw trainable
spin-matrix parameters into a Clifford-structured spin operator. Abstractly,
one may write
\[
  \mathcal{C}(W)
  =
  \sum_{a} c_a(W)\,\Gamma_a,
\]
where $\{\Gamma_a\}$ is a Clifford basis and $c_a(W)$ are complex coefficients
encoded by the learnable matrix. This is useful when the learned spin operation
should be expressed in a basis natural for Dirac spinors rather than as an
unstructured dense matrix.

\section{Core convolution-like layers}

The layers in this section are the most direct analogue of convolutional
layers. Their receptive field is specified by paths on the lattice, but the
field values collected from those paths are made comparable by gauge parallel
transport. This is the central technical distinction between \LMP{} and an
ordinary convolutional network.

\subsection[PTC Layer: parallel-transport convolution]{\texttt{PTC\_Layer}: parallel-transport convolution}

\paragraph{Implementation.}
Defined in \code{src/Layer_types.jl}; forward map in
\code{src/Layer_sequence.jl}; functional helper \code{PTC} in
\code{src/elementry_base.jl}.

\paragraph{Inputs and parameters.}
Let $\paths=\{p_1,\dots,p_K\}$ be the paths stored in a \code{Hopping_operator}.
Let $W_i\in\C^{4\times4}$ be trainable spin matrices. The layer supports either
one shared matrix $W$ or one matrix per path.

\paragraph{Forward map.}
The path-transported features are
\[
  h_i = \pathop_{p_i}[\U]\field.
\]
The layer computes
\[
  y
  =
  \sigma\left(
  \sum_{i=1}^{K} W_i h_i
  \right),
\]
or
\[
  y
  =
  \sigma\left(
  \sum_{i=1}^{K} W h_i
  \right)
\]
when a single shared weight is used. If \code{residual=true}, the output is
\[
  y \leftarrow y+\field.
\]

\paragraph{Scientific idea.}
This is the LQCD analogue of a convolution. Ordinary convolution translates
neighboring pixels to a common location and mixes them. In a gauge theory, a
fermion value at a neighboring site cannot be directly compared with a fermion
at the current site, because the two live in different local color frames.
\code{PTC_Layer} first parallel-transports each neighbor to the target site,
then learns a spin-space mixing. The result is a local, gauge-covariant
feature extractor.

\subsection[CPTC Layer: Clifford parallel-transport convolution]{\texttt{CPTC\_Layer}: Clifford parallel-transport convolution}

\paragraph{Implementation.}
Defined in \code{src/Layer_types.jl}; forward map in
\code{src/Layer_sequence.jl}; helper \code{CPTC} in
\code{src/elementry_base.jl}.

\paragraph{Forward map.}
\code{CPTC_Layer} has the same transport structure as \code{PTC_Layer}, but it
uses Clifford-transformed weights:
\[
  y
  =
  \sigma\left(
  \sum_{i=1}^{K} \mathcal{C}(W_i)\,\pathop_{p_i}[\U]\field
  \right).
\]

\paragraph{Scientific idea.}
Wilson fermions carry spinor degrees of freedom, and the Wilson-Dirac operator
itself is built from gamma matrices. A Clifford-structured layer biases the
learned spin operation toward this physical algebra. It is a physics-informed
restriction of the spin-space part of the convolution.

\subsection[PTLayer: multi-input path transport layer]{\texttt{PTLayer}: multi-input path transport layer}

\paragraph{Implementation.}
Defined in \code{src/Layer_types.jl}; forward map in
\code{src/Layer_sequence.jl}; CPU and GPU paths are both implemented.

\paragraph{Inputs and parameters.}
\code{PTLayer} acts on a vector of fermion fields
\[
  x=(x_1,\dots,x_K),
\]
where the number of inputs matches the number of paths. It has an output count
$O$ and trainable weights $W_{j i}\in\C^{4\times4}$.

\paragraph{Forward map.}
For each input $x_i$ and path $p_i$, compute
\[
  h_i = \pathop_{p_i}[\U]x_i.
\]
Then for output channel $j$,
\[
  y_j
  =
  \sum_{i=1}^{K}
  W_{j i} h_i,
  \qquad j=1,\dots,O.
\]
The implementation stores weights in a flattened order, but the mathematical
map is a channel-mixing operator after gauge-covariant transport.

\paragraph{Scientific idea.}
\code{PTLayer} generalizes \code{PTC_Layer} from a single input field to a
multi-channel fermion representation. It is useful when a model carries several
candidate low modes, coarse vectors, or learned feature fields through a shared
gauge background.

\subsection[DenseLinearLayer: spin-matrix dense channel mixing]{\texttt{DenseLinearLayer}: spin-matrix dense channel mixing}

\paragraph{Implementation.}
Defined in \code{src/Layer_types.jl}; forward map in
\code{src/Layer_sequence.jl}.

\paragraph{Forward map.}
For input fields $x_1,\dots,x_I$ and output fields $y_1,\dots,y_O$,
\[
  y_j
  =
  \sum_{k=1}^{I}
  W_{j k} x_k,
  \qquad W_{j k}\in\C^{4\times4}.
\]
There is no gauge-link transport in this layer; it assumes all fields already
transform at the same lattice site.

\paragraph{Scientific idea.}
This is a dense linear layer in feature-channel space, with spin matrices as
the learned coefficients. Because it only mixes spin and channel indices, it
preserves gauge covariance of each field channel.

\section{Attention-style layers}

The attention layers should not be read as a direct copy of transformer
attention on token embeddings. In \LMP{}, Q/K/V features are fermion fields
or spin matrices derived from fermion fields. Their contractions must be chosen
so that the resulting weights are meaningful under a change of local gauge
frame. The attention mechanism is therefore best understood as a learned
selection among gauge-transported paths and spin channels.

\subsection{General covariant attention principle}

Attention in ordinary machine learning builds scores from query and key
features, then uses those scores to weight value features. In LQCD, all of
these features must transform coherently under local color rotations. \LMP{}
therefore constructs Q/K/V features by applying gauge-covariant hopping and
spin-space weights. Typical features have the form
\[
  Q_i = W_i^Q\,\pathop_{p_i}[\U]\field,\quad
  K_i = W_i^K\,\pathop_{p_i}[\U]\field,\quad
  V_i = W_i^V\,\pathop_{p_i}[\U]\field.
\]
Gauge-invariant or spin-local scores are then made from contractions such as
\[
  M_i = \frac{1}{m}\,\mathcal{M}(Q_i^\dagger,K_i),
\]
where $m$ is the layer's \code{nominator} scaling factor and $\mathcal{M}$ is
implemented by \code{Matmul!} or \code{spincolortrace}, depending on the layer.

\subsection[Attention Layer]{\texttt{Attention\_Layer}}

\paragraph{Implementation.}
Defined in \code{src/Layer_types.jl}; forward map starts at
\code{Attention_Layer} in \code{src/Layer_sequence.jl}.

\paragraph{Forward map.}
The code first Clifford-transforms the Q/K/V weights:
\[
  \widetilde W_i^Q=\mathcal{C}(W_i^Q),\quad
  \widetilde W_i^K=\mathcal{C}(W_i^K),\quad
  \widetilde W_i^V=\mathcal{C}(W_i^V).
\]
For each path feature $h_i=\pathop_{p_i}[\U]\field$,
\[
  Q_i=\widetilde W_i^Qh_i,\qquad
  K_i=\widetilde W_i^Kh_i.
\]
The score is
\[
  S_i
  =
  \sigma\left(
  \frac{\operatorname{Matmul}(Q_i^\dagger,K_i)}{m}
  \right),
\]
then it is multiplied by the value spin weight and applied to $h_i$:
\[
  y
  =
  \sum_i
  \left(S_i\widetilde W_i^V\right) h_i.
\]
If residual mode is enabled, $y\leftarrow y+\field$.

\paragraph{Scientific idea.}
This layer learns a spin-structured importance matrix for each transported
path. It can be read as a local self-attention over covariantly transported
fermion features, where the attention weight is itself a spin matrix rather
than a scalar probability.

\subsection[FullPTCAttention Layer]{\texttt{FullPTCAttention\_Layer}}

\paragraph{Forward map.}
\code{FullPTCAttention_Layer} is similar in spirit but uses the raw PTC weights
for Q/K/V construction:
\[
  Q_i = W_i^Q\,\pathop_{p_i}[\U]\field,\qquad
  K_i = W_i^K\,\pathop_{p_i}[\U]\field.
\]
The score and output are
\[
  S_i = \sigma\left(
  \frac{\operatorname{Matmul}(Q_i^\dagger,K_i)}{m}
  \right),
  \qquad
  y = \sum_i (S_i W_i^V)\,\pathop_{p_i}[\U]\field.
\]

\paragraph{Scientific idea.}
This is closer to a direct Q/K/V attention mechanism over path-transported
features. It provides more freedom than the Clifford-restricted variant.

\subsection[PrallelTransportAttention Layer]{\texttt{PrallelTransportAttention\_Layer}}

\paragraph{Name note.}
The public type name currently uses the spelling \code{PrallelTransport...} in
the source and exports. This document keeps the implementation spelling.

\paragraph{Forward map.}
The implementation builds
\[
  K_i = W_i^K\,\pathop_{p_i}[\U]\field,\qquad
  h_i = \pathop_{p_i}[\U]\field.
\]
Using the identity-path or first hopped feature $h_1$ as a reference, it forms
\[
  M_i =
  \frac{\operatorname{Matmul}(h_1^\dagger,K_i)}{m}.
\]
It extracts diagonal magnitudes
\[
  R_{i,\alpha}= |(M_i)_{\alpha\alpha}|,
  \qquad
  r_\alpha=\sum_i R_{i,\alpha},
\]
then applies a thresholded ratio with bias $b$:
\[
  S_{i,\alpha\alpha}
  =
  \max\left(\frac{R_{i,\alpha}}{r_\alpha}-b,0\right)
  (M_i)_{\alpha\alpha}.
\]
The output is
\[
  y = \sum_i (S_i W_i^V)h_i,
\]
with optional residual addition.

\paragraph{Scientific idea.}
This layer selects transported paths according to their relative diagonal
spin-channel response. The bias parameter acts as a sparsifying threshold,
encouraging only paths with strong relative contribution to enter the value
sum.

\subsection[siteweight Layer]{\texttt{siteweight\_Layer}}

\paragraph{Forward map.}
This layer computes Q/K/V path features and uses a local color-spin contraction:
\[
  S_{pq}(n)
  =
  \frac{1}{m}
  \operatorname{spincolortrace}
  \left(Q_p^\dagger(n),K_q(n)\right).
\]
The result is assembled as
\[
  y(n)
  =
  \sum_{p,q}
  S_{pq}(n)\,V_q(n).
\]

\paragraph{Scientific idea.}
Unlike the spin-matrix attention variants, \code{siteweight_Layer} produces
site-dependent scalar or tensor weights from local contractions. This is closer
to spatial attention over the lattice while retaining gauge-compatible
contractions.

\section{Spectral and variational layers}

The spectral layers turn a neural network into a variational ansatz for
Dirac-operator eigenspaces. Instead of comparing the output only to a fixed
target field, these layers measure whether the output fields span a physically
useful low-mode subspace. This is the part of the package most closely tied to
deflation, multigrid setup, and solver acceleration.

\subsection{Rayleigh quotient motivation}

Many LQCD workflows depend on low modes of Dirac-type operators because small
eigenvalues dominate solver convergence, multigrid setup quality, and
deflation efficiency. \LMP{} includes layers that evaluate Rayleigh-type
objectives to train fermion fields or feature channels toward low-mode
subspaces.

For a linear operator $\dirac$, a Rayleigh quotient for the positive operator
$\dirac^\dagger\dirac$ is
\[
  R(x)
  =
  \frac{\inner{\dirac x}{\dirac x}}{\inner{x}{x}}
  =
  \frac{\inner{x}{\dirac^\dagger\dirac x}}{\inner{x}{x}}.
\]
Minimizing this quantity encourages $x$ to align with low singular/eigenmodes
of the Dirac operator.

\subsection[RayleighQuotient Layer]{\texttt{RayleighQuotient\_Layer}}

\paragraph{Implementation.}
Defined in \code{src/Layer_types.jl}; forward map in
\code{src/Layer_sequence.jl}.

\paragraph{Forward map.}
For the stored field $x$ and input operator $\dirac$,
\[
  R(x;\dirac)
  =
  \frac{\norm{\dirac x}^{2}}{\norm{x}^{2}}.
\]

\subsection[PtcRayleighQuotient Layer]{\texttt{PtcRayleighQuotient\_Layer}}

\paragraph{Forward map.}
For a vector of fields $x_1,\dots,x_K$ and stored operator $\dirac$,
\[
  L(x_1,\dots,x_K)
  =
  \sum_{i=1}^{K}
  \frac{\norm{\dirac x_i}^{2}}{\norm{x_i}^{2}}.
\]
This lets a multi-channel network optimize several candidate vectors.

\subsection[HoppingRayleighQuotient Layer]{\texttt{HoppingRayleighQuotient\_Layer}}

\paragraph{Forward map.}
This layer first creates a sum of hopped fields
\[
  y = \sum_{i=1}^{K}\pathop_{p_i}[\U]x,
\]
then evaluates
\[
  R(y;\dirac)
  =
  \frac{\norm{\dirac y}^{2}}{\norm{y}^{2}}.
\]

\paragraph{Scientific idea.}
The quotient is evaluated after a gauge-covariant path expansion. This allows
the model to assess whether transported local combinations improve alignment
with low modes.

\subsection[Ritz Layer]{\texttt{Ritz\_Layer}}

\paragraph{Implementation.}
Defined in \code{src/Layer_types.jl}; forward map in
\code{src/Layer_sequence.jl}; GPU code uses CUSOLVER for the generalized
Hermitian eigenproblem.

\paragraph{Forward map.}
Given a set of candidate vectors $V=\{x_1,\dots,x_l\}$, the layer builds Gram
matrices
\begin{align}
  A_{ij} &= \inner{\dirac x_i}{\dirac x_j},\\
  B_{ij} &= \inner{x_i}{x_j}.
\end{align}
The Ritz values are obtained from the generalized eigenproblem
\[
  A c_k = \lambda_k B c_k.
\]
The layer returns the real sum
\[
  L_{\mathrm{Ritz}}(V)
  =
  \sum_{k=1}^{l}\lambda_k.
\]

\paragraph{Scientific idea.}
The Ritz procedure optimizes an entire subspace rather than isolated vectors.
For LQCD, this is useful for learning low-mode subspaces suitable for
deflation, coarse-grid construction, or solver preconditioning.

\section{Orthogonalization and normalization}

The layers in this section are not primarily feature extractors. They enforce
or encourage numerical conditions that are important when learning several
fermion fields simultaneously. Without such structure, multiple trainable
vectors can collapse onto the same low mode or develop inconvenient scales.

\subsection[GramSchmidt Layer]{\texttt{GramSchmidt\_Layer}}

\paragraph{Forward map.}
The layer performs modified Gram-Schmidt orthogonalization:
\begin{align}
  q_1 &= x_1,\\
  q_i &=
  x_i -
  \sum_{j<i}
  \frac{\inner{q_j}{x_i}}{\inner{q_j}{q_j}}q_j,
  \qquad i=2,\dots,l.
\end{align}
It returns $(q_1,\dots,q_l)$.

\paragraph{Scientific idea.}
Ritz and Rayleigh objectives can collapse multiple vectors toward the same
mode. Orthogonalization keeps a set of learned vectors linearly independent,
which is required for stable subspace learning.

\subsection[Orthogonal Layer]{\texttt{Orthogonal\_Layer}}

\paragraph{Forward map.}
For coefficient $a$, this scalar objective is
\[
  L_{\mathrm{orth}}(x_1,\dots,x_l)
  =
  a
  \sum_{1\le j<i\le l}
  \left|\inner{x_j}{x_i}\right|^2.
\]

\paragraph{Scientific idea.}
This is a soft penalty that discourages non-orthogonal vectors without directly
replacing them by an orthogonalized basis.

\subsection[NormLayer]{\texttt{NormLayer}}

\paragraph{Forward map.}
For each input field $x_i$ and scale $\alpha$,
\[
  y_i = \alpha \frac{x_i}{\sqrt{\norm{x_i}^{2}}}.
\]

\paragraph{Scientific idea.}
Rayleigh-type losses are scale invariant, and uncontrolled vector norms can
hurt numerical stability. \code{NormLayer} explicitly fixes the field norm.

\subsection{Other normalization layers}

\begin{itemize}[leftmargin=*]
  \item \code{complexNormLayer}: subtracts the field mean,
        $x\mapsto x-\operatorname{mean}(x)$.
  \item \code{zscoreNormLayer}: applies standardization.
  \item \code{minmaxNormLayer}: applies min-max scaling,
        $x\mapsto (x-\min x)/(\max x-\min x)$.
\end{itemize}

\section{Block and domain-decomposition layer}

\subsection[BlockScalar Layer]{\texttt{BlockScalar\_Layer}}

\paragraph{Implementation.}
Defined in \code{src/Layer_types.jl}; forward map in
\code{src/Layer_sequence.jl}; block utilities in \code{src/elementry_base.jl}.

\paragraph{Forward map.}
The lattice is decomposed into blocks of linear size $r$. Let $b(n)$ be the
block index containing site $n$. For each input field $x_k$, the layer has a
complex scalar weight $w_{b,k}$ for every block. The essential operation is
\[
  y_k(n)
  =
  w_{b(n),k}\,x_k(n).
\]
When residual mode is enabled, the input field is added back to the output.

\paragraph{Scientific idea.}
Blockwise scalar weights are a simple local coarse representation. They are
related to domain decomposition and multigrid ideas: local regions of the
lattice can receive separate amplitudes while preserving the internal
color-spin structure of the fermion field. This is useful for learning local
coarse-vector corrections or block-adaptive preconditioner components.

\section{Dirac and clover-related operators}

\subsection[Dirac Layer]{\texttt{Dirac\_Layer}}

\code{Dirac_Layer} wraps an operator and applies it to a fermion field:
\[
  y = \dirac x.
\]
This is useful when a model needs to insert a known physics operator as a
non-trainable layer.

\subsection{Clover Dirac support}

\code{src/elementry_base.jl} also contains utilities for a clover-improved
Dirac operator. The documented formula in the source is
\[
  D_{\mathrm{cl}}
  =
  D_{\mathrm{Wilson}}
  +
  i c_{\mathrm{SW}}\kappa
  \sum_{\mu<\nu}
  \sigma_{\mu\nu}\widehat F_{\mu\nu},
  \qquad
  \sigma_{\mu\nu}
  =
  \frac{i}{2}[\gamma_\mu,\gamma_\nu],
\]
with
\[
  \widehat F_{\mu\nu}
  =
  \frac{Q_{\mu\nu}-Q_{\mu\nu}^{\dagger}}{8}.
\]
Although this is not a trainable layer by itself, it is part of the physics
operator infrastructure that can appear inside spectral objectives such as
\code{Ritz_Layer}.

\section{Activation functions}

The layer structs store an \code{ActivationFunction}. The currently visible
types include:
\begin{itemize}[leftmargin=*]
  \item \code{Identity}: returns the input.
  \item \code{Softmax}: row-wise exponential normalization for spin matrices.
  \item \code{ModRelu}: applies a modulus threshold to complex spin-matrix
        entries,
        \[
          z \mapsto \max(|z|-0.1,0)\,e^{i\arg z}.
        \]
\end{itemize}
These activations mostly act on spin matrices or spin-matrix collections in the
attention layers, while many field-valued layers default to identity.

\section{Layer summary table}

\begin{table}[h]
\centering
\small
\begin{tabular}{@{}p{0.27\linewidth}p{0.39\linewidth}p{0.25\linewidth}@{}}
\toprule
Layer & Main mathematical object & Primary role \\
\midrule
\code{PTC_Layer} & $\sum_i W_i\pathop_{p_i}\field$ & Gauge-covariant convolution \\
\code{CPTC_Layer} & $\sum_i \mathcal{C}(W_i)\pathop_{p_i}\field$ & Clifford-informed convolution \\
\code{PTLayer} & $y_j=\sum_i W_{ji}\pathop_{p_i}x_i$ & Multi-channel path transport \\
\code{DenseLinearLayer} & $y_j=\sum_k W_{jk}x_k$ & Spin/channel mixing \\
\code{Attention_Layer} & $S_i=\sigma(Q_i^\dagger K_i/m)$ & Covariant spin attention \\
\code{FullPTCAttention_Layer} & Raw Q/K/V PTC attention & Flexible path attention \\
\code{PrallelTransportAttention_Layer} & Diagonal ratio-threshold scores & Sparse path selection \\
\code{siteweight_Layer} & $\operatorname{spincolortrace}(Q^\dagger,K)$ & Site-dependent weighting \\
\code{RayleighQuotient_Layer} & $\norm{\dirac x}^2/\norm{x}^2$ & Low-mode objective \\
\code{Ritz_Layer} & $A c=\lambda B c$ & Subspace low-mode objective \\
\code{GramSchmidt_Layer} & Modified Gram-Schmidt & Basis orthogonalization \\
\code{BlockScalar_Layer} & $w_{b(n),k}x_k(n)$ & Block/domain scaling \\
\bottomrule
\end{tabular}
\caption{Main layer families in \LMP{} and their scientific roles.}
\end{table}

\section{Implementation map}

\begin{itemize}[leftmargin=*]
  \item Public exports: \code{src/LMP.jl}.
  \item Layer type definitions: \code{src/Layer_types.jl}.
  \item Forward maps and layer sequencing: \code{src/Layer_sequence.jl}.
  \item Path representation and hopping: \code{src/elementry_base.jl}.
  \item Gauge transformation: \code{src/Gauge.jl}.
  \item GPU helpers and batched kernels: \code{src/gpu.jl}.
  \item Basic attention tests: \code{test/Attention_layer.jl} and
        \code{test/Attention_gradient.jl}.
  \item Hopping covariance test: \code{test/hopping_covariance_test.jl}.
\end{itemize}

\section{Practical interpretation}

The most important conceptual point is that \LMP{} layers are not generic
neural-network layers applied to arrays. They are physics-aware operators on
fermion fields. The package repeatedly uses the following pattern:
\[
  \begin{aligned}
  &\text{transport by gauge links}\\
  &\quad\longrightarrow\ \text{spin/channel mixing}\\
  &\quad\longrightarrow\ \text{spectral or variational objective}.
  \end{aligned}
\]
This pattern makes the learned maps compatible with local gauge covariance and
connects the learned representation to LQCD solver tasks, especially low-mode
learning, deflation, and multigrid-inspired coarse representations.

\section{Suggested narrative for collaborators}

When introducing \LMP{} to another LQCD researcher, a concise description is:
\begin{quote}
\LMP{} explores neural-network layers whose elementary operation is not a
standard array convolution, but a gauge-covariant transport-and-mix operation
on Wilson fermion fields. A layer first gathers fields along gauge-link paths,
parallel-transports them to a common color frame, and then applies learnable
spin/channel matrices. On top of this feature construction, the package defines
attention-like path selection and Rayleigh--Ritz objectives for learning
low-mode subspaces of Dirac-type operators.
\end{quote}

This framing emphasizes that the package is not simply using machine learning
as a black box. The learned maps are constrained by the same geometric
structure that appears in LQCD operators. Consequently, the output fields and
objectives have a direct interpretation in terms of gauge-covariant features,
spinor algebra, and spectral approximation.

\section{Current status and cautions}

The document describes the implemented layer families as they appear in the
current source tree. A few points are worth keeping in mind when using the
document as research material:
\begin{itemize}[leftmargin=*]
  \item Several layers expose both CPU and GPU paths. The mathematical maps are
        the same, but the implementation may use cached buffers or fused kernels
        for performance.
  \item The spelling \code{PrallelTransportAttention_Layer} is preserved because
        it is the public name currently used in the source.
  \item Some activation functions are implemented for spin matrices rather than
        general fermion fields. In many field-valued layers the default
        activation is identity.
  \item The formulas here describe the forward maps and scientific intent. They
        are not a full derivation of every gradient implementation.
\end{itemize}

\section{Concluding perspective}

\LMP{} is best understood as an attempt to make neural architectures speak the
language of lattice gauge theory. Its layers are built from parallel transport,
spinor mixing, local contractions, and variational spectral objectives. These
are precisely the structures that appear in Wilson-Dirac operators, deflation
methods, and multigrid algorithms. For that reason, the package is a promising
framework for testing whether trainable gauge-covariant maps can produce useful
low-mode or coarse-space information for LQCD computations.

\end{document}
